% Chapter 5: Testing and Results

This chapter presents the testing methodologies and experimental evaluation performed on the Hybrid Post-Quantum VPN system. The testing process validates the correctness of the hybrid cryptographic handshake, encrypted tunnel communication, Linux TUN interface operation, frontend-backend integration, and overall system reliability. The experimental results demonstrate the feasibility of integrating post-quantum cryptography into practical VPN deployments while maintaining acceptable performance overhead.

\section{Cryptographic Testing}
\label{sec:crypto_testing}

Cryptographic testing validates the correctness of the hybrid handshake implementation, key derivation process, encryption pipeline, and authentication mechanisms.

\subsection{Test Cases}

\begin{table}[H]
\centering
\begin{tabularx}{\textwidth}{|X|X|l|}
\hline
\textbf{Test Description} & \textbf{Expected Result} & \textbf{Status} \\ \hline
Generate X25519 key pairs & Valid public/private keys generated & Pass \\ \hline
Perform ML-KEM-768 encapsulation & Shared secret and ciphertext generated & Pass \\ \hline
Perform ML-KEM-768 decapsulation & Client and server shared secrets match & Pass \\ \hline
Combine classical and PQ secrets using HKDF & Identical session keys derived on both nodes & Pass \\ \hline
ECDSA-P256 signature verification & Handshake transcript verified successfully & Pass \\ \hline
AES-256-GCM encryption/decryption & Packet decrypted correctly without corruption & Pass \\ \hline
Replay packet with duplicate nonce & Packet rejected by replay protection & Pass \\ \hline
Tamper encrypted packet authentication tag & Packet integrity verification failed & Pass \\ \hline
\end{tabularx}
\caption{Cryptographic Test Cases}
\end{table}

\subsubsection{Sample Test: Hybrid Handshake Validation}

\begin{lstlisting}[language=Python, caption={Test: Hybrid Handshake Validation}]
# Generate classical shared secret
x25519_secret = perform_x25519_exchange()

# Generate post-quantum shared secret
mlkem_secret = perform_mlkem_encapsulation()

# Derive final session key
session_key = hkdf_sha256(
    x25519_secret + mlkem_secret
)

assert client_session_key == server_session_key

# Result: PASS
\end{lstlisting}

\begin{lstlisting}[language=Python, caption={Test: AES-256-GCM Integrity Validation}]
# Encrypt packet
ciphertext, nonce, tag = aes_gcm_encrypt(packet)

# Tamper authentication tag
tag = b"invalidtag"

# Attempt decryption
decrypt_packet(ciphertext, nonce, tag)

# Result: Authentication failed, packet dropped - PASS
\end{lstlisting}

\section{Tunnel and Networking Testing}
\label{sec:tunnel_testing}

Tunnel testing validates Linux TUN interface creation, encrypted UDP communication, routing functionality, and packet forwarding.

\subsection{Test Cases}

\subsubsection{TUN Interface Testing}

\begin{table}[H]
\centering
\begin{tabularx}{\textwidth}{|X|X|l|}
\hline
\textbf{Test Description} & \textbf{Expected Result} & \textbf{Status} \\ \hline
Create Linux TUN interface & Interface \texttt{hyb0} created successfully & Pass \\ \hline
Assign IP address to TUN interface & Correct IP assigned using pyroute2 & Pass \\ \hline
Configure routing table & Tunnel traffic routed through TUN device & Pass \\ \hline
Destroy tunnel session & TUN interface cleaned up correctly & Pass \\ \hline
\end{tabularx}
\caption{TUN Interface Test Cases}
\end{table}

\subsubsection{UDP Tunnel Testing}

\begin{table}[H]
\centering
\begin{tabularx}{\textwidth}{|X|X|l|}
\hline
\textbf{Test Description} & \textbf{Expected Result} & \textbf{Status} \\ \hline
Transmit encrypted UDP packets & Packets received successfully at gateway & Pass \\ \hline
Forward packets through tunnel & Bidirectional communication established & Pass \\ \hline
Packet replay attack simulation & Duplicate packets rejected & Pass \\ \hline
Invalid nonce transmission & Packet discarded immediately & Pass \\ \hline
Gateway disconnect detection & Tunnel terminated within timeout period & Pass \\ \hline
\end{tabularx}
\caption{UDP Tunnel Test Cases}
\end{table}

\subsubsection{Performance Testing}

\begin{table}[H]
\centering
\begin{tabularx}{\textwidth}{|X|X|l|}
\hline
\textbf{Performance Metric} & \textbf{Observed Result} & \textbf{Status} \\ \hline
Hybrid handshake latency overhead & Less than 50 ms additional latency & Pass \\ \hline
AES-256-GCM throughput & Sustained throughput above 50 Mbps & Pass \\ \hline
Average packet encryption latency & Stable under nominal load & Pass \\ \hline
Tunnel establishment time & Completed successfully within expected duration & Pass \\ \hline
\end{tabularx}
\caption{Performance Testing Results}
\end{table}

\section{Frontend and API Testing}
\label{sec:frontend_testing}

Frontend testing validates the Electron + Next.js dashboard, API communication, user interactions, and tunnel management controls.

\subsection{Test Cases}

\subsubsection{Dashboard Testing}

\begin{table}[H]
\centering
\begin{tabularx}{\textwidth}{|X|X|l|}
\hline
\textbf{Test Description} & \textbf{Expected Result} & \textbf{Status} \\ \hline
Load Electron dashboard & Dashboard rendered successfully & Pass \\ \hline
Display tunnel connection status & Real-time status updated correctly & Pass \\ \hline
Connect to gateway & Handshake initiated successfully & Pass \\ \hline
Disconnect active tunnel & Tunnel terminated cleanly & Pass \\ \hline
Display cryptographic suite & Active algorithms shown correctly & Pass \\ \hline
Display byte counters and metrics & Real-time metrics updated correctly & Pass \\ \hline
\end{tabularx}
\caption{Dashboard Test Cases}
\end{table}

\subsubsection{API Testing}

\begin{table}[H]
\centering
\begin{tabularx}{\textwidth}{|X|X|l|}
\hline
\textbf{Test Description} & \textbf{Expected Result} & \textbf{Status} \\ \hline
Call \texttt{/health} endpoint & Service health returned successfully & Pass \\ \hline
Call \texttt{/status} endpoint & Tunnel state returned correctly & Pass \\ \hline
Initiate demo handshake & Handshake completed successfully & Pass \\ \hline
Connect to gateway API & Session established successfully & Pass \\ \hline
Disconnect from tunnel & Tunnel shutdown completed successfully & Pass \\ \hline
\end{tabularx}
\caption{API Endpoint Test Cases}
\end{table}

\section{System Testing}
\label{sec:system_testing}

System testing validates the complete end-to-end workflow of the Hybrid Post-Quantum VPN prototype across multiple virtual machines.

\subsection{End-to-End Tunnel Establishment}

This test validates the complete VPN workflow between the client VM and gateway VM.

\begin{enumerate}
    \item \textbf{Gateway VM:}
    \begin{itemize}
        \item Starts the Gateway API service.
        \item Initializes UDP listener and TUN interface.
    \end{itemize}

    \item \textbf{Client VM:}
    \begin{itemize}
        \item Starts the Agent API and Electron dashboard.
        \item Initiates hybrid handshake with the gateway.
    \end{itemize}

    \item \textbf{Handshake Phase:}
    \begin{itemize}
        \item X25519 and ML-KEM-768 operations executed successfully.
        \item Session keys derived using HKDF-SHA256.
        \item ECDSA authentication verified.
    \end{itemize}

    \item \textbf{Tunnel Phase:}
    \begin{itemize}
        \item Linux TUN interfaces created successfully.
        \item IPv4 packets encrypted and transmitted over UDP.
        \item AES-256-GCM integrity checks validated packets successfully.
    \end{itemize}
\end{enumerate}

\textbf{Result:} Secure encrypted communication established successfully between both virtual machines using hybrid cryptography. \textbf{PASS}

\subsection{Cross-Platform Testing}

\begin{table}[H]
\centering
\begin{tabularx}{\textwidth}{|X|X|l|}
\hline
\textbf{Test Description} & \textbf{Expected Result} & \textbf{Status} \\ \hline
Run backend on Ubuntu Linux & Full functionality available & Pass \\ \hline
Run frontend in browser mode & Dashboard functions correctly & Pass \\ \hline
Run backend on Windows/macOS & Handshake functions without TUN support & Pass \\ \hline
Docker container deployment & APIs accessible successfully & Pass \\ \hline
\end{tabularx}
\caption{Cross-Platform Test Cases}
\end{table}

\subsection{Security Testing}

\begin{table}[H]
\centering
\begin{tabularx}{\textwidth}{|X|X|l|}
\hline
\textbf{Test Description} & \textbf{Expected Result} & \textbf{Status} \\ \hline
Replay attack simulation & Duplicate packets rejected & Pass \\ \hline
Handshake tampering attempt & Transcript verification failed & Pass \\ \hline
Modify encrypted payload & AES-GCM authentication failed & Pass \\ \hline
Invalid signature verification & Connection rejected & Pass \\ \hline
Nonce reuse simulation & Rekeying triggered successfully & Pass \\ \hline
\end{tabularx}
\caption{Security Test Cases}
\end{table}

\section{Experimental Results}
\label{sec:results}

The experimental evaluation demonstrates that hybrid cryptographic integration can be achieved with acceptable overhead while significantly improving resistance against future quantum threats.

\subsection{Observed Results}

\begin{enumerate}
    \item The hybrid handshake successfully combined X25519 and ML-KEM-768 shared secrets using HKDF-SHA256.

    \item AES-256-GCM encryption provided authenticated encrypted communication with stable throughput.

    \item Replay protection mechanisms correctly rejected duplicate or tampered packets.

    \item The Linux TUN-based encrypted tunnel operated successfully between virtual machines.

    \item The Electron + Next.js dashboard successfully monitored and controlled tunnel operations.

    \item Docker deployment enabled simplified testing and reproducibility of the backend services.
\end{enumerate}

\subsection{Comparison with Classical VPNs}

\begin{table}[H]
\centering
\begin{tabular}{|l|c|c|}
\hline
\textbf{Feature} & \textbf{Classical VPN} & \textbf{Hybrid PQC VPN} \\ \hline
Quantum Resistance & No & Yes \\ \hline
Classical Security & Yes & Yes \\ \hline
Handshake Complexity & Lower & Moderate \\ \hline
Forward Secrecy & Yes & Yes \\ \hline
Post-Quantum Protection & No & Yes \\ \hline
Algorithm Agility & Limited & High \\ \hline
\end{tabular}
\caption{Comparison Between Classical and Hybrid VPNs}
\end{table}

\section{Test Summary}

\begin{table}[H]
\centering
\begin{tabular}{|l|c|c|}
\hline
\textbf{Test Category} & \textbf{Total Tests} & \textbf{Pass Rate} \\ \hline
Cryptographic Testing & 8 & 100\% \\ \hline
Tunnel and Networking Testing & 14 & 100\% \\ \hline
Frontend and API Testing & 11 & 100\% \\ \hline
System and Security Testing & 10 & 100\% \\ \hline
\textbf{Total} & \textbf{43} & \textbf{100\%} \\ \hline
\end{tabular}
\caption{Overall Test Execution Summary}
\end{table}

All testing phases completed successfully, demonstrating that the Hybrid Post-Quantum VPN prototype provides secure encrypted communication, functional hybrid cryptographic integration, and reliable tunnel operation while supporting future migration toward quantum-safe networking infrastructures.